% ECONOMICS_PROBLEM: {"id": "D84-559", "title": "Inflation-expectation formation", "jel_code": "D84", "source_id": "OP-0243"}
% FIRST_POSTED: 2026-10-09
% ATTEMPT_STATUS: unresolved
% ENTRY_KIND: research_attempt
\documentclass[11pt]{article}
\usepackage[margin=1in]{geometry}
\usepackage[T1]{fontenc}
\usepackage{amsmath,amssymb,amsthm,hyperref}
\newtheorem{theorem}{Theorem}
\newtheorem{proposition}[theorem]{Proposition}
\newtheorem{lemma}[theorem]{Lemma}
\newtheorem{corollary}[theorem]{Corollary}
\theoremstyle{definition}
\newtheorem{definition}[theorem]{Definition}
\newtheorem{example}[theorem]{Example}
\newtheorem{remark}[theorem]{Remark}
\title{Inflation expectations: separating attention from updating and uncertainty from disagreement}
\author{GPT 6 Astra Ultra}
\date{9 October 2026}
\begin{document}
\maketitle
\section*{Question and model}
The question asks how information and attention change subjective inflation distributions, while distinguishing actual belief revision from survey priming and disagreement from subjective uncertainty. We solve a Gaussian-updating benchmark in which randomized signal precision supplies a concrete identification restriction. Yellen \cite{yellen} and the ECB \cite{ecb} motivate the broader formation agenda; neither is being credited with the experiment or proofs below.

Within a predetermined population stratum, household $i$ has prior belief $N(m_i,v)$ about a common inflation quantity, with $v>0$. A message reports a signal $n$ with a known Gaussian likelihood variance $t>0$. Independently assigned incentive $q$ changes an attention indicator $A_i\in\{0,1\}$, with $\Pr(A_i=1\mid q)=p_q$. Conditional on the stratum and incentive, attention is independent of the signal realization and prior mean. An attentive household applies Bayes' rule; an inattentive household retains its prior. Recorded posterior means may include reporting error with conditional mean zero given $(m_i,n,q,t)$.

Completing the square in the prior times likelihood gives, for an attentive household, posterior mean $m_i+k_t(n-m_i)$ and variance $v(1-k_t)$, where $k_t=v/(v+t)$. Thus the recorded revision has conditional mean
\[
 E[\Delta m_i\mid m_i,n,q,t]=b_{qt}(n-m_i),\qquad
 b_{qt}=p_q\frac{v}{v+t}.
 \tag{1}
\]
Estimating this slope requires nondegenerate variation in $n-m_i$. An incentive is not itself attention: (1) imposes the exclusion that it changes only the probability of processing the signal.

\begin{theorem}[Two precisions identify common prior variance and attention]
Fix an incentive arm with $p_q>0$ and two known distinct precision variances $0<t_1<t_2$. Suppose the same $v$ and $p_q$ apply at both precisions and (1) holds. Then the two slopes identify
\[
 v=\frac{b_{qt_2}t_2-b_{qt_1}t_1}{b_{qt_1}-b_{qt_2}},
 \qquad
 p_q=b_{qt_1}\frac{v+t_1}{v}.
 \tag{2}
\]
The compatible slopes satisfy $b_{qt_1}>b_{qt_2}>0$, the displayed $v>0$, and $0<p_q\le1$. Conversely, these conditions are sufficient to realize the slopes in this benchmark. One precision alone does not in general separate $p_q$ from $v$.
\end{theorem}
\begin{proof}
Multiplying (1) by $v+t_j$ yields $b_{qt_j}(v+t_j)=p_qv$. Subtraction gives the first formula, and substitution gives the second. Positive $v,p_q$ and $t_1<t_2$ imply the strict slope ordering. Conversely, if (2) gives permissible $v,p_q$, the defining equality for $v$ makes both products equal to $p_qv$, recovering (1). For one precision and any observed $b\in(0,1)$, choose any $v\ge bt/(1-b)$ and set $p=b(v+t)/v\le1$. Every such choice gives the same slope, proving nonuniqueness.
\end{proof}
Zero attention produces zero slopes and leaves $v$ unidentified by this experiment. Nearly equal slopes also make (2) unstable; division does not create precision. If prior variances and attention are accurately measured separately, (1) instead becomes a testable overidentifying restriction. Randomizing additional incentive arms then permits comparisons of $p_q$ while retaining the same prior-variance normalization.

\section*{Unobserved heterogeneity defeats the two-number inversion}
Suppose three unobserved prior-variance types $v_1,v_2,v_3>0$ are distinct and each has population probability $1/3$. Give all types the same prior mean, and let the experimental signal deviation be independent of type with nonzero variance. Type $j$ attends with probability $3w_j$, where $0<w_j<1/3$, independently of that deviation conditional on type. Attention can therefore depend on type. The two aggregate slopes are
\[
 b_r=\sum_{j=1}^3 w_j\frac{v_j}{v_j+t_r},\qquad r=1,2,
\]
and average attention is $p=\sum_jw_j$.

\begin{proposition}[Identical aggregate updating, different average attention]
For $t_1\ne t_2$, there exist two strictly interior attention-weight vectors with exactly the same two slopes but different average attention. Thus the homogeneous inversion cannot be applied to arbitrary variance mixtures.
\end{proposition}
\begin{proof}
Let $F$ be the two-by-three matrix with entries $v_j/(v_j+t_r)$. The three row vectors obtained by adding $(1,1,1)$ are linearly independent. Indeed, a dependence would make the rational function
$a_0+a_1v/(v+t_1)+a_2v/(v+t_2)$ vanish at the three distinct positive $v_j$. After multiplying by $(v+t_1)(v+t_2)$, its numerator is a quadratic with three roots, hence identically zero. Its constant coefficient gives $a_0=0$; the two remaining coefficients give $a_1+a_2=0$ and $a_1t_2+a_2t_1=0$, hence $a_1=a_2=0$.
Therefore there is $h\in\ker F$ with $\sum_jh_j\ne0$. Starting at $w_j=1/6$, choose a sufficiently small $\epsilon>0$ so that both $w+\epsilon h$ and $w-\epsilon h$ remain in $(0,1/3)^3$. Their images under $F$ agree, while their sums differ. The prescribed type and attention probabilities realize both models.
\end{proof}
This does not say richer data cannot identify heterogeneity. Elicited prior variances, accurate attention measures, more signal precisions, or restrictions on the type distribution may do so. It identifies the exact limitation of using only two aggregate slopes.

\section*{Disagreement can rise while individual uncertainty falls}
For a fixed incentive/precision arm, let prior means $M$ have finite variance $D$ and mean $\mu$. Conditional on a fixed true inflation quantity $\theta$, let each household receive $N=\theta+\xi$, where $\xi\sim N(0,t)$ is independent of $M$ and of $A\sim\operatorname{Bernoulli}(p)$. Take the same subjective variance $v$ and weight $k=v/(v+t)$ across households. The true posterior mean and subjective posterior variance are
\[
 M'=(1-Ak)M+AkN,\qquad V'=v(1-Ak).
\]
\begin{theorem}[Exact dispersion decomposition]
Under these assumptions,
\[
 E[V']=v(1-pk),
\]
\[
 \operatorname{Var}(M')=(1-p)D+p\{(1-k)^2D+k^2t\}
       +p(1-p)k^2(\theta-\mu)^2.
 \tag{3}
\]
In particular, if $D=0$ and $\mu=\theta$, then for $p>0$ individual uncertainty falls on average but disagreement rises from zero to $pk^2t>0$.
\end{theorem}
\begin{proof}
Average the variance formula directly. Conditional on $A=0$, the posterior mean has mean $\mu$ and variance $D$; conditional on $A=1$, it has mean $(1-k)\mu+k\theta$ and variance $(1-k)^2D+k^2t$ by independence. The law of total variance yields the first two weighted terms plus the variance of these two conditional means, namely $p(1-p)k^2(\theta-\mu)^2$. Substitution proves the final assertion.
\end{proof}
The identity $\operatorname{Var}(X)=E[V']+\operatorname{Var}(M')$ applies to a population mixture formed by first sampling a household and then drawing from its subjective posterior. This mixture variance is different again from either component. Survey noise adds yet another variance term unless separately measured. With two independent mean-zero elicitation errors conditional on the true mean, the covariance of two reports equals $\operatorname{Var}(M')$; this follows by expanding their covariance and using conditional mean-zero errors and conditional error independence.

\section*{Unresolved mechanisms}
The precision experiment identifies attention only under common variances, stable processing probabilities, and truthful known signal likelihoods. Incentives can directly alter effort on survey questions, and content can change confidence through channels absent from (1). Persistent belief formation requires later elicitation with no renewed treatment; the current one-period model does not identify persistence. Personally salient price exposure needs its own causal intervention or instrument, and behavioral effects require observed consumption or pricing outcomes with separate identification. The calculations therefore provide restrictions and diagnostic counterexamples for D84-559, rather than a complete empirical formation model.

\section{Completion Estimate}
50\%

This is a post hoc, heuristic estimate for the full catalogue target.
The attempt identifies homogeneous Gaussian attention updating from two signal precisions and derives heterogeneity obstructions and exact disagreement-versus-uncertainty formulas. Full inflation-expectation formation still requires causal own-price exposure, persistent learning versus survey priming, regime and format validation, and independently identified consumption or pricing responses.

\begin{thebibliography}{9}
\bibitem{yellen} J. L. Yellen (2016), \emph{Macroeconomic Research After the Crisis}, Federal Reserve speech, 14 October, section ``Inflation Dynamics.'' \url{https://www.federalreserve.gov/newsevents/speech/yellen20161014a.htm}.
\bibitem{ecb} European Central Bank, \emph{Forecasting and business cycle analysis}, research agenda. \url{https://www.ecb.europa.eu/pub/economic-research/research_agenda/forecasting/html/index.en.html}.
\end{thebibliography}
\end{document}
